\section{The algebraic determinant and its arithmetic bound}
\label{sec:arithmetic}

Fix \(\nu>2\kappa\) and assume \eqref{eq:bad-approximations}.
For now take finite parameters \(m,K,w_0,v_0,\theta\) satisfying the
interpolation hypotheses, and choose good approximants
\(p_i/q_i\), \(q_i\ge2\), \(p_i\ne0\), with separated
\(w_i=\lceil\log q_i\rceil\). Their precise order of choice is
given in Section~\ref{sec:parameters}.
Let \(F_0>2/\theta\), and set
\begin{equation}\label{eq:truncations}
 T_i=\left\lceil\frac{F_0w_i}{v_0}\right\rceil,\qquad
 G_i(t)=\sum_{1\le k<T_i}\frac{(-1)^{k+1}}{k}t^k,\qquad
 L_i=\lcm(1,\ldots,T_i).
\end{equation}
Using \(T_i\), rather than \(T_i-1\), in the least common multiple
is harmless and agrees with the formal arithmetic bound.
The omitted tail has weight at least \(v_0T_i\ge F_0w_i>w_i/\theta\).

\subsection{Matrix, rank, and one minor over the number field}

Index columns by \((h,\alpha)\in\N\times\N^m\) satisfying
\(w_0h+w\alpha\le H\), and rows by
\((j,\ell,\beta)\) satisfying
\[
 0\le j<K,\qquad v_0\ell+w\beta/\theta<H.
\]
Define the matrix over \(F\) by
\begin{equation}\label{eq:matrix}
 M_H[(j,\ell,\beta),(h,\alpha)]
 =[t^\ell u^\beta]\,
   \bigl(\xi^j(1+t)\bigr)^h
   \prod_{i=1}^m
    \bigl(j\gamma p_i/q_i+G_i(t)+u_i\bigr)^{\alpha_i}.
\end{equation}
Explicitly, for \(\beta\le\alpha\), its entry is
\begin{equation}\label{eq:binomial-entry}
 \xi^{jh}\binom{\alpha}{\beta}
 [t^\ell](1+t)^h
   \prod_i\bigl(j\gamma p_i/q_i+G_i(t)\bigr)^{\alpha_i-\beta_i};
\end{equation}
it is zero otherwise.

Fix one embedding in \(S\). If the powers of \(\xi\) are distinct,
the centers
\[
 \bigl(\sigma(\xi)^j,
       (j\sigma(\gamma)p_i/q_i)_{i=1}^m\bigr)
\]
satisfy case (i) of Theorem~\ref{thm:interpolation}.
If \(\xi=1\), each additive coordinate is injective in \(j\)
because \(\gamma p_i/q_i\ne0\), and case (ii) applies.
The truncation in \eqref{eq:truncations} preserves the jet quotient.
It follows that, for sufficiently large divisible \(H\), the
complex image of \eqref{eq:matrix} has full row rank.
Select a square minor using every row, and denote its determinant
by \(\Delta_H\in F^\times\) and its size by \(M=M(H)\).
The chosen columns may depend on \(H\).
Every embedding sends this same determinant to the determinant of
the corresponding complex minor; none of these images is zero.
We never choose different minors at different embeddings.

The fixed-weight simplex counts give
\begin{equation}\label{eq:row-count}
 M(H)\sim
 \frac{K\theta^m H^{m+1}}{(m+1)!v_0\prod_iw_i}.
\end{equation}
Define
\begin{equation}\label{eq:mean}
 \bar b_H=\frac1{MH}\sum_{(j,\ell,\beta)\ {\rm row}}w\beta,
 \qquad 0\le\bar b_H\le\theta.
\end{equation}
The simplex asymptotic follows, for example, by comparing unit cubes
with the bounded shifts of the corresponding real simplex. Changing
a weak boundary inequality to a strict one does not change its leading
volume.

\subsection{Clearing the base, slope, and logarithmic denominators}

Choose positive integers \(D_\xi,D_\gamma\) such that
\(D_\xi\xi,D_\gamma\gamma\in\OO_F\), and put
\begin{equation}\label{eq:common-denominator}
 \mathcal D_H=
 D_\xi^{\,K\lfloor H/w_0\rfloor}
 \prod_{i=1}^m(D_\gamma L_i)^{\lfloor H/w_i\rfloor}.
\end{equation}
Scale each column by \(\prod_iq_i^{\alpha_i}\), each row by
\(\prod_iq_i^{-\beta_i}\), and every entry by \(\mathcal D_H\).

\begin{lemma}\label{lem:integral}
Every scaled entry belongs to \(\OO_F\). Hence the scaled determinant
is a nonzero algebraic integer
\[
 Z_H=S_H\Delta_H,\qquad
 S_H=\mathcal D_H^M
   \prod_{\text{selected columns}}\prod_iq_i^{\alpha_i}
   \prod_{\text{rows}}\prod_iq_i^{-\beta_i}>0.
\]
Here \(S_H\in\Q\), and it need not be an integer.
\end{lemma}
\begin{proof}
Entries with \(\beta\not\le\alpha\) vanish.
For the others, the row and column factors turn
\eqref{eq:binomial-entry} into
\[
 \xi^{jh}\binom{\alpha}{\beta}
 [t^\ell](1+t)^h
   \prod_i(j\gamma p_i+q_iG_i(t))^{\alpha_i-\beta_i}.
\]
The polynomial
\(D_\gamma L_i(j\gamma p_i+q_iG_i(t))\)
has integral coefficients. Its required power is at most
\(\alpha_i\le\lfloor H/w_i\rfloor\).
Also \(jh\le K\lfloor H/w_0\rfloor\), so
\[
 D_\xi^{K\lfloor H/w_0\rfloor}\xi^{jh}
 =(D_\xi\xi)^{jh}D_\xi^{K\lfloor H/w_0\rfloor-jh}
\]
is integral. This proves entrywise integrality.
The determinant of an integral matrix is integral, and the scalings
are nonzero.
\end{proof}

Let \(\Lambda=\log4+4\). We use the elementary bound
\(\log\lcm(1,\ldots,n)\le\Lambda n\).
Indeed, the elementary Chebyshev bound
\(\vartheta(u)\le u\log4\), obtained from binomial coefficients,
and the prime-power identity for \(\psi\) give, for \(n\ge1\),
\[
 \log\lcm(1,\ldots,n)=\psi(n)
 \le n\log4+2\sqrt n\log n
 \le(\log4+4)n.
\]
The last step uses \(\log n\le2\sqrt n\); \(n=0\) is immediate.
This deliberately loose constant is exactly
\path{OAI.PiExponent.Arithmetic.lcmConstant} in the pinned source.
The sharper constant in the informal argument of
\cite[Section 3.1]{OAIpi} is not needed.
Since \(T_i\le F_0w_i/v_0+1\),
\begin{equation}\label{eq:log-D}
 \frac{\log\mathcal D_H}{H}
 \le \frac{\Lambda F_0m}{v_0}
    +(\Lambda+\log D_\gamma)\sum_i\frac1{w_i}
    +\frac{K\log D_\xi}{w_0}.
\end{equation}
Write \(w_*=\min_iw_i\). For every selected column,
\(\sum_i\alpha_i\log q_i\le H\).
For a row, \(\log q_i\ge w_i-1\) and
\(|\beta|\le\theta H/w_*\), so
\[
 \sum_i\beta_i\log q_i\ge w\beta-\theta H/w_*.
\]
Consequently
\begin{equation}\label{eq:scaling-bound}
 \frac{\log S_H}{MH}\le1-\bar b_H+E_{\rm ar}^{\,0},
 \qquad
 E_{\rm ar}^{\,0}=
 \frac{\Lambda F_0m}{v_0}
 +(\Lambda+\log D_\gamma)\sum_i\frac1{w_i}
 +\frac{K\log D_\xi}{w_0}+\frac{\theta}{w_*}.
\end{equation}

\begin{proposition}[Norm lower bound]\label{prop:arithmetic}
Set
\[
 \mathcal N_H=\frac1{sMH}
        \sum_{\sigma:F\hookrightarrow\C}\log|\sigma(\Delta_H)|,
 \qquad E_{\rm ar}=\kappa E_{\rm ar}^{\,0}.
\]
Then
\begin{equation}\label{eq:arithmetic}
 \mathcal N_H\ge-\kappa(1-\bar b_H)-E_{\rm ar}.
\end{equation}
\end{proposition}
\begin{proof}
The norm of the nonzero algebraic integer \(Z_H\) is a nonzero
rational integer. Therefore
\[
 1\le|N_{F/\Q}(Z_H)|
   =S_H^d\prod_{\sigma:F\hookrightarrow\C}|\sigma(\Delta_H)|.
\]
Take logarithms, divide by \(sMH\), and apply
\eqref{eq:scaling-bound}.
\end{proof}

For \(x=\log3\), take \(F=\Q\), \(\gamma=1\), \(\xi=3\),
and \(D_\xi=D_\gamma=1\): the factor \(3^{jh}\) is integral
and costs no denominator. For a positive rational \(r=a/b\),
one can take \(D_\xi=b\); the extra cost is at most
\(K\log b/w_0\), which the parameter choice will make small.
